\documentclass[12pt]{article}

\usepackage{fontawesome}
\usepackage{hyperref}
\usepackage[tagged, highstructure]{accessibility}
\usepackage{bookmark}
\usepackage{enumitem}

\hypersetup{
    colorlinks=false,
    pdfborder={0 0 0},
    pdftitle={Installs, Updates, and Documentation Lab},
    pdfauthor={Adrianna Holden-Gouveia},
    pdfsubject={Linux Administration - Lab Assignment},
    pdflang={en-US},
    pdfkeywords={Linux installs, updates, package management, Fail2Ban, repositories}

}

% Configure PDF bookmarks for navigation
\bookmarksetup{
    numbered,
    open,
}

% Configure list spacing for better accessibility
\setlist{nosep}



\title{Installs, Updates, and Documentation}
\author{
        Adrianna Holden-Gouveia \\
        Website: \url{https://aholdengouveia.name}\\ 
        \date{\vspace{-5ex}}
        %Email: \href{mailto:admin@aholdengouveia.name}{admin@aholdengouveia.name} \\
%        \faLinkedin{: aholdengouveia} \\
        \faGithub {: aholdengouveia} \\
        %\faTwitter {: aholdengouveia} \\
        }

%S\date{\today}


\begin{document}    

\maketitle

%\begin{abstract}
%This is a template for Linux Administration Lab work
%\end{abstract}
%\tableofcontents


\section*{Accessibility Notice}
This document is also available in HTML format at:

\url{https://aholdengouveia.name/LinuxAdmin/labexcercises/contSetup.html}

The HTML version provides enhanced accessibility features including keyboard navigation, screen reader support, responsive design, dark mode support, and high contrast options.

\section*{Objectives:}
\begin{itemize}
    \item The objective of this lab is to introduce students to essential Linux system administration tasks including package management, software installation, and system updates. Through hands-on exercises across multiple servers, students will learn to manage repositories, install and configure various software packages, and maintain comprehensive documentation of system changes—critical skills for effective server administration and team collaboration.
\end{itemize}
\section*{Complete the following problems}

References, a video, a PowerPoint and some notes are available at my website
\url{https://www.aholdengouveia.name/LinuxAdmin/installs.html}



\subsection*{Go through both servers and complete the following problems.  Make sure your deliverables include documents for BOTH servers clearly labeled with what the document is and which server it's for}
    \begin{enumerate}
        \item Check your repositories that are included by default on both servers, find out how you'd go about adding a repository, (both official and PPA) and make sure you can install programs.
        \item Install Tmux, emacs, and Fail2Ban.  Make sure you look through the Fail2Ban settings as you do the install, and note anything that needs to be changed or modified to get Fail2Ban up and working.
        \item Just for fun, install cowsay and lolcat.   (Yes I know these are silly, the point is that i know they aren't there already.)
        \item Update VI to vim, update your servers if it's needed.
        \item Documentation should include locations of any files you've created or need to reference. You should include any programs you've installed or updated as well, with dates of when you've done those things. You also need to make sure to show how to open each program, what they look like when open, how to use them, and how to close each program.
    \end{enumerate}

\section*{Deliverables}
Target Market is someone trying to do the installs on their servers as well, assume technical knowledge but not expertise.  Screenshots are helpful to go with your descriptions. 

\begin{itemize}
    \item A document showing how you installed everything on both servers.
    \item A short document showing how to find where your repos live on both servers as well as how to add/modify a repo or PPA.
    \item A short document outlining what you've installed on your servers,  including the programs listed here. Make sure to show how to find everything installed on the server.
    \item A short document including any updates and dependencies that needed to be done or changed and the dates that you did this.  Anything you've turned on/off should be included as well, don't forget the dates that you did this
    \item For Fail2ban make sure you note what you've changed (if anything) and why you made that change.
\end{itemize}


\end{document}